% Part: methods
% Chapter: induction
% Section: relations

\documentclass[../../../include/open-logic-section]{subfiles}

\begin{document}

\olfileid[hi]{mth}{ind}{rel}

\olsection{संबंध और फलन}

जब हमने वस्तुओं के किसी समुच्चय (जैसे प्राकृतिक संख्याएँ या सुंदर पद)
को आगमनात्मक रूप से परिभाषित कर लिया हो, तो इन वस्तुओं \emph{पर संबंध} भी
आगमन से परिभाषित कर सकते हैं। उदाहरणतः इस विचार पर विचार करें: सुंदर पद
$t_1$, सुंदर पद $t_2$ का उपपद है यदि वह उसके एक भाग के रूप में आता है। इसके
लिए चिह्न उपयोग करें: $t_1\sqsubseteq t_2$। निस्संदेह प्रत्येक सुंदर पद
स्वयं का उपपद है: $t\sqsubseteq t$। इस संबंध की आगमनात्मक परिभाषा इस प्रकार
दे सकते हैं:

\begin{defn}
  सुंदर पद $t_1$ के $t_2$ का उपपद होने का संबंध $t_1\sqsubseteq t_2$,
  $t_2$ पर आगमन से इस प्रकार परिभाषित है:
  \begin{enumerate}
  \item यदि $t_2$ कोई अक्षर है, तो $t_1\sqsubseteq t_2$ तभी जब $t_1=t_2$।
  \item यदि $t_2$, $[s_1\circ s_2]$ है, तो $t_1\sqsubseteq t_2$ तभी जब
    $t_1=t_2$, $t_1\sqsubseteq s_1$, अथवा $t_1\sqsubseteq s_2$।
  \end{enumerate}
\end{defn}

उदाहरणतः यह परिभाषा हमें बताएगी कि
$\mathrm{a}\sqsubseteq[\mathrm{b}\circ\mathrm{a}]$। क्योंकि (2) कहता है
कि $\mathrm{a}\sqsubseteq[\mathrm{b}\circ\mathrm{a}]$ तभी जब
$\mathrm{a}=[\mathrm{b}\circ\mathrm{a}]$, अथवा
$\mathrm{a}\sqsubseteq\mathrm{b}$, अथवा
$\mathrm{a}\sqsubseteq\mathrm{a}$। पहले दो असत्य हैं: स्पष्टतः
$\mathrm{a}$, $[\mathrm{b}\circ\mathrm{a}]$ के समान नहीं है, और (1) से
$\mathrm{a}\sqsubseteq\mathrm{b}$ तभी जब $\mathrm{a}=\mathrm{b}$, जो भी
असत्य है। किंतु (1) से ही $\mathrm{a}\sqsubseteq\mathrm{a}$ तभी जब
$\mathrm{a}=\mathrm{a}$, जो सत्य है।

यह ध्यान देना महत्त्वपूर्ण है कि इस परिभाषा की सफलता एक ऐसे तथ्य पर निर्भर
है जिसे हमने अभी सिद्ध नहीं किया: प्रत्येक सुंदर पद $t$ या तो स्वयं एक
अक्षर है, अथवा ऐसे \emph{अद्वितीय रूप से निर्धारित} सुंदर पद $s_1$ और $s_2$
हैं कि $t=[s_1\circ s_2]$। यहाँ ``अद्वितीय रूप से निर्धारित'' का अर्थ है
कि यदि $t=[s_1\circ s_2]$, तो $s_1\neq r_1$ या $s_2\neq r_2$ वाले
$[r_1\circ r_2]$ के \emph{भी} बराबर नहीं। यदि ऐसा होता, तो खंड (2) स्वयं
से टकरा सकता था: $t_2$ को $[s_1\circ s_2]$ पढ़कर हमें
$t_1\sqsubseteq t_2$ मिल सकता, किंतु $t_2$ को $[r_1\circ r_2]$ पढ़कर
$t_1\not\sqsubseteq t_2$ मिल सकता। यह सिद्ध करने से पहले कि ऐसा नहीं हो
सकता, एक ऐसा उदाहरण देखें जहाँ ऐसा \emph{हो सकता है}।

\begin{defn}
  \emph{कोष्ठक-रहित पदों} को आगमनात्मक रूप से इस प्रकार परिभाषित करें
  \begin{enumerate}
  \item प्रत्येक अक्षर एक कोष्ठक-रहित पद है।
    \item यदि $s_1$ और $s_2$ कोष्ठक-रहित पद हैं, तो $s_1\circ s_2$ एक
      कोष्ठक-रहित पद है।
    \item और कुछ भी कोष्ठक-रहित पद नहीं है।
  \end{enumerate}
\end{defn}

कोष्ठक-रहित पदों के उदाहरण हैं $\mathrm{a}$,
$\mathrm{b}\circ\mathrm{d}$, $\mathrm{b}\circ\mathrm{a}\circ\mathrm{b}$।
अब यदि कोष्ठक-रहित पदों के लिए ``उपपद'' को ऊपर की तरह परिभाषित करें, तो
दूसरा खंड होगा
\begin{center}
  यदि $t_2=s_1\circ s_2$, तो $t_1\sqsubseteq t_2$ तभी जब $t_1=t_2$,
  $t_1\sqsubseteq s_1$, अथवा $t_1\sqsubseteq s_2$।
\end{center}

अब $\mathrm{b}\circ\mathrm{a}\circ\mathrm{b}$ इस प्रकार $s_1\circ s_2$
रूप का है
\begin{align*}
  s_1 & = \mathrm{b} \text{ और} & s_2 & = \mathrm{a} \circ \mathrm{b}.
\intertext{यह इस प्रकार $r_1\circ r_2$ रूप का भी है}
  r_1 & = \mathrm{b} \circ \mathrm{a} \text{ और} & r_2 &= \mathrm{b}.
\end{align*}
अब क्या $\mathrm{a}\circ\mathrm{b}$,
$\mathrm{b}\circ\mathrm{a}\circ\mathrm{b}$ का उपपद है? पहले पाठ के अनुसार
उत्तर हाँ है और दूसरे के अनुसार नहीं।

जिस गुण के अनुसार अन्य सुंदर पदों से किसी सुंदर पद के निर्माण का ढंग
अद्वितीय होता है, उसे \emph{अद्वितीय पठनीयता} कहते हैं। चूँकि ऐसी
आगमनात्मक रूप से परिभाषित वस्तुओं के संबंधों की आगमनात्मक परिभाषाएँ
महत्त्वपूर्ण हैं, हमें सिद्ध करना होगा कि यह गुण सत्य है।

\begin{prop}
  मान लें $t$ एक सुंदर पद है। तब या तो $t$ स्वयं एक अक्षर है, अथवा ऐसे
  अद्वितीय रूप से निर्धारित सुंदर पद $s_1,s_2$ हैं कि
  $t=[s_1\circ s_2]$।
\end{prop}

\begin{proof}
  यदि $t$ स्वयं एक अक्षर है, तो शर्त संतुष्ट है। अतः मान लें $t$ स्वयं एक
  अक्षर नहीं। आगमनात्मक परिभाषा से तब $t$ किन्हीं सुंदर पदों $s_1$ और $s_2$
  के लिए $[s_1\circ s_2]$ रूप का होना चाहिए। यह दिखाना शेष है कि ये
  अद्वितीय रूप से निर्धारित हैं, अर्थात् यदि $t=[r_1\circ r_2]$, तो
  $s_1=r_1$ और $s_2=r_2$।

  अतः सुंदर पदों $s_1,s_2,r_1,r_2$ के लिए मान लें
  $t=[s_1\circ s_2]$ और $t=[r_1\circ r_2]$ भी। हमें दिखाना है कि
  $s_1=r_1$ और $s_2=r_2$। पहले $s_1$ और $r_1$ समान होने चाहिए, क्योंकि
  अन्यथा एक दूसरे का उचित आरंभिक खंड होगा। किंतु
  \olref[sti]{prop:initial} से यह असंभव है यदि $s_1$ और $r_1$ दोनों सुंदर
  पद हों। और यदि $s_1=r_1$, तो स्पष्टतः $s_2=r_2$ भी।
\end{proof}

हम फलनों को भी आगमनात्मक रूप से परिभाषित कर सकते हैं: उदाहरणतः वह फलन $f$
जो किसी सुंदर पद को उसमें अंतर्निहित $[\dots]$ की अधिकतम गहराई पर
प्रतिचित्रित करता है, इस प्रकार परिभाषित कर सकते हैं:

\begin{defn}
  \ollabel{defn:depth} सुंदर पद की \emph{गहराई} $f(t)$ आगमनात्मक रूप से इस
  प्रकार परिभाषित है:
  \[
  f(t) = \begin{cases}
    0 & \text{ यदि $t$ एक अक्षर है}\\
    \max(f(s_1), f(s_2)) + 1 & \text{ यदि $t=[s_1\circ s_2]$।}
  \end{cases}
  \]
\end{defn}

उदाहरणतः
\begin{align*}
  f([\mathrm{a} \circ \mathrm{b}]) & = 
    \max(f(\mathrm{a}),f(\mathrm{b})) + 1 = \\ 
  &= \max(0, 0) + 1 = 1, \text{ और}\\
  f([[\mathrm{a} \circ \mathrm{b}] \circ \mathrm{c}]) & = 
    \max(f([\mathrm{a} \circ \mathrm{b}]), f(\mathrm{c})) + 1 = \\ 
  & = \max(1,0) + 1 = 2.
\end{align*}

यहाँ निस्संदेह हम मानते हैं कि $s_1$ और $s_2$ सुंदर पद हैं और इस तथ्य का
उपयोग करते हैं कि प्रत्येक सुंदर पद या तो अक्षर है अथवा
$[s_1\circ s_2]$ रूप का। फिर यह महत्त्वपूर्ण है कि वह इस रूप का केवल एक
ढंग से हो सकता है। क्यों, यह देखने के लिए पहले परिभाषित कोष्ठक-रहित पदों पर
फिर विचार करें। संगत ``परिभाषा'' होगी:
\[
  g(t) = 
  \begin{cases}
    0 & \text{ यदि $t$ एक अक्षर है}\\
   \max(g(s_1), g(s_2)) + 1 & \text{ यदि $t=s_1\circ s_2$।}
  \end{cases}
\]
अब कोष्ठक-रहित पद
$\mathrm{a}\circ\mathrm{b}\circ\mathrm{c}\circ\mathrm{d}$ पर विचार करें।
उसे एक से अधिक ढंग से पढ़ा जा सकता है, उदाहरणतः इस प्रकार $s_1\circ s_2$
\begin{align*}
  s_1 & = \mathrm{a} \text{ और} & 
  s_2 & = \mathrm{b} \circ \mathrm{c} \circ \mathrm{d},
\intertext{अथवा इस प्रकार $r_1\circ r_2$}
  r_1 & = \mathrm{a} \circ b \text{ और} & 
  r_2 &= \mathrm{c} \circ \mathrm{d}.
\end{align*}
पहले पाठ के अनुसार $g$ का परिकलन देगा
\begin{align*}
  g(s_1 \circ s_2) & =
  \max(g(\mathrm{a}), g(\mathrm{b} \circ \mathrm{c} \circ \mathrm{d})) + 1 =\\ & =
  \max(0,2) + 1 = 3
  \intertext{जबकि दूसरे पाठ के अनुसार हमें मिलता है}
  g(r_1 \circ r_2) & =
  \max(g(\mathrm{a} \circ \mathrm{b}), g(\mathrm{c} \circ \mathrm{d})) + 1=\\ &
  = \max(1,1) + 1 = 2
\end{align*}
किंतु किसी फलन को सदैव अद्वितीय मान देना चाहिए; अतः $g$ की हमारी
``परिभाषा'' किसी फलन को बिल्कुल परिभाषित नहीं करती।

\begin{prob}
  फलन $l$ की आगमनात्मक परिभाषा दीजिए, जहाँ $l(t)$ सुंदर पद $t$ में चिह्नों
  की संख्या है।
\end{prob}

\begin{prob}
  सुंदर पदों $t$ पर संरचनात्मक आगमन से सिद्ध कीजिए कि $f(t)<l(t)$ (जहाँ
  $l(t)$, $t$ में चिह्नों की संख्या है और $f(t)$,
  \olref[mth][ind][rel]{defn:depth} में परिभाषित $t$ की गहराई है)।
\end{prob}

\end{document}
